\documentclass{article}

\usepackage[english]{babel}
\usepackage[letterpaper,top=2cm,bottom=2cm,left=3cm,right=3cm,marginparwidth=1.75cm]{geometry}
\setlength{\parindent}{0pt}

\usepackage{amsmath}
\usepackage{amsthm}
\usepackage{amssymb}
\usepackage{graphicx}
\usepackage{float}
\usepackage{listings}
\usepackage[colorlinks=true, allcolors=blue]{hyperref}
\usepackage{xcolor}

\lstset{
  basicstyle=\ttfamily\small,
  numbers=left,
  numberstyle=\tiny,
  stepnumber=1,
  numbersep=8pt,
  backgroundcolor=\color{gray!10},
  frame=single,
  rulecolor=\color{gray!60},
  tabsize=2,
  captionpos=b,
  breaklines=true,
  breakatwhitespace=true,
  showstringspaces=false,
  keywordstyle=\color{blue},
  commentstyle=\color{green!50!black},
  stringstyle=\color{orange},
}

\newtheorem{theorem}{Theorem}[section]
\newtheorem{lemma}[theorem]{Lemma}
\newtheorem{proposition}[theorem]{Proposition}
\newtheorem{corollary}[theorem]{Corollary}
\theoremstyle{definition}
\newtheorem{definition}[theorem]{Definition}
\theoremstyle{remark}
\newtheorem{remark}[theorem]{Remark}

\newcommand{\EE}{\mathbb{E}}
\newcommand{\PP}{\mathbb{P}}
\newcommand{\Var}{\mathrm{Var}}
\newcommand{\cH}{\mathcal{H}}
\newcommand{\cD}{\mathcal{D}}
\newcommand{\ind}{\mathbf{1}}

\title{DSC 291: Learning Theory Assignment 3}
\author{Tongle Shen}

\begin{document}
\maketitle

\section{Part A: A Mini-Course on Concentration Inequalities}

\subsection{Notation and Setup}

Throughout this mini-course I use one notation system across all four inequalities.

\begin{itemize}
\item $X_1,\dots,X_n$ are random variables.
\item $\mu_i = \EE[X_i]$ and, when all means are the same, $\mu = \EE[X_i]$.
\item The centered sum is
$$
S_n = \sum_{i=1}^n (X_i - \mu_i).
$$
\item The empirical average is
$$
\overline{X}_n = \frac{1}{n} \sum_{i=1}^n X_i.
$$
\item For Bernstein's inequality, the total variance proxy is
$$
V_n = \sum_{i=1}^n \Var(X_i).
$$
\item For bounded-difference inequalities, $f(X_1,\dots,X_n)$ is a function of independent inputs, and $c_i$ denotes the sensitivity of $f$ to the $i$th coordinate.
\end{itemize}

When I switch to learning-theoretic notation later, I will do so explicitly and then reduce back to these same concentration templates.

\subsection{The Concept of Concentration}

\subsubsection{What Concentration Means}

A concentration inequality is a statement of the form
$$
\PP(|Y - \EE[Y]| \ge t) \le \text{small tail bound in } t,
$$
where $Y$ is a random quantity built from many smaller random ingredients.

The point is not merely that $Y$ has an expectation. The point is that $Y$ is unlikely to wander far from that expectation. When the tail bound decays exponentially in $t^2$ or in $t$, the mean becomes a genuinely useful summary of the random variable.

\subsubsection{Why This Should Happen}

There are two recurring reasons concentration appears.

\textbf{Averaging effect.} If $Y$ is a sum or average of many independent or weakly dependent terms, positive and negative fluctuations tend to cancel.

\textbf{Stability effect.} If $Y=f(X_1,\dots,X_n)$ and changing one input changes $f$ only a little, then no single coordinate can move $Y$ very far. This is the viewpoint behind McDiarmid's inequality.

\subsubsection{Concrete Example}

Let $X_1,\dots,X_n$ be i.i.d.\ Bernoulli$(p)$ variables, so each $X_i \in \{0,1\}$ and $\EE[X_i]=p$. Then
$$
\overline{X}_n = \frac{1}{n}\sum_{i=1}^n X_i
$$
is the sample frequency of success, while $p$ is the true success probability.

We expect $\overline{X}_n$ to be close to $p$ because each individual flip contains limited information and the average pools many such observations. Hoeffding's inequality makes that intuition quantitative:
$$
\PP(|\overline{X}_n-p| \ge \varepsilon) \le 2e^{-2n\varepsilon^2}.
$$
So the failure probability decays exponentially fast in $n$.

\subsection{The Common Technique: MGF and Chernoff Bounds}

\subsubsection{The MGF Point of View}

For a real-valued random variable $Y$, its moment generating function is
$$
M_Y(\lambda) = \EE[e^{\lambda Y}],
\qquad \lambda \in \mathbb{R}
$$
whenever this expectation exists.

The MGF is useful because large values of $Y$ are amplified exponentially by $e^{\lambda Y}$. If we can show that $\EE[e^{\lambda Y}]$ is small enough, then large positive excursions of $Y$ must be rare.

\subsubsection{Chernoff Bound as the General Template}

\begin{lemma}[Chernoff bound template]
\label{lem:chernoff-template}
Let $Y$ be any random variable for which $\EE[e^{\lambda Y}] < \infty$ for some $\lambda > 0$. Then for every $t>0$,
$$
\PP(Y \ge t) \le e^{-\lambda t}\EE[e^{\lambda Y}].
$$
Consequently,
$$
\PP(Y \ge t) \le \inf_{\lambda > 0} e^{-\lambda t}\EE[e^{\lambda Y}].
$$
\end{lemma}

\begin{proof}
By Markov's inequality applied to the nonnegative random variable $e^{\lambda Y}$,
$$
\PP(Y \ge t)
=
\PP(e^{\lambda Y} \ge e^{\lambda t})
\le
e^{-\lambda t}\EE[e^{\lambda Y}].
$$
Now optimize over $\lambda > 0$.
\end{proof}

So every MGF-based proof has the same skeleton:
\begin{enumerate}
\item choose the random quantity $Y$ whose tail you want to control;
\item bound $\EE[e^{\lambda Y}]$;
\item plug that bound into Lemma~\ref{lem:chernoff-template};
\item optimize in $\lambda$.
\end{enumerate}

\subsubsection{What Needs to Be Controlled}

To get a useful tail bound, we want the centered MGF
$$
\EE[e^{\lambda (Y-\EE[Y])}]
$$
to look like the MGF of a well-understood distribution.

\begin{itemize}
\item If it is bounded by $e^{c\lambda^2}$, then we obtain Gaussian-type tails $e^{-t^2/(4c)}$.
\item If it is bounded by something like $e^{\lambda^2 v /(2(1-\lambda b/3))}$, then we obtain Bernstein-type tails, which use both a variance term $v$ and a range bound $b$.
\end{itemize}

\subsubsection{The Bounded-Differences Variant}

For McDiarmid's inequality the random quantity is not usually a plain sum. Instead we write
$$
f(X_1,\dots,X_n)-\EE[f(X_1,\dots,X_n)]
$$
as a martingale sum
$$
\sum_{i=1}^n D_i
$$
where each increment $D_i$ has conditional mean zero and controlled range. Then we run the same MGF method conditionally, one increment at a time.

So the real common template is broader than ``sum of independent variables.'' It is:

\begin{quote}
Represent the fluctuation as a sum of increments with controlled exponential moments, and then apply Chernoff.
\end{quote}

\subsection{Theorem Statements, Assumptions, and Consequences}

\subsubsection{Hoeffding's Inequality}

\begin{theorem}[Hoeffding inequality \cite{Hoeffding1963}]
\label{thm:hoeffding}
Let $X_1,\dots,X_n$ be independent random variables such that
$$
a_i \le X_i \le b_i
\qquad \text{almost surely for } i=1,\dots,n.
$$
Then for every $t>0$,
$$
\PP\!\left(
\sum_{i=1}^n (X_i-\mu_i) \ge t
\right)
\le
\exp\!\left(
-\frac{2t^2}{\sum_{i=1}^n (b_i-a_i)^2}
\right).
$$
Equivalently, for every $\varepsilon > 0$,
$$
\PP\!\left(
\overline{X}_n - \frac{1}{n}\sum_{i=1}^n \mu_i \ge \varepsilon
\right)
\le
\exp\!\left(
-\frac{2n^2\varepsilon^2}{\sum_{i=1}^n (b_i-a_i)^2}
\right).
$$
The same bound holds for the lower tail, and therefore the two-sided bound is
$$
\PP\!\left(
\left|
\sum_{i=1}^n (X_i-\mu_i)
\right| \ge t
\right)
\le
2\exp\!\left(
-\frac{2t^2}{\sum_{i=1}^n (b_i-a_i)^2}
\right).
$$
\end{theorem}

This theorem controls sums of \emph{independent bounded variables}. It does not use any variance information beyond boundedness, so it is robust but sometimes conservative.

\subsubsection{Sampling Without Replacement}

\begin{theorem}[Hoeffding inequality for sampling without replacement \cite{Hoeffding1963,Serfling1974}]
\label{thm:hoeffding-without-replacement}
Let $z_1,\dots,z_N \in [0,1]$ be fixed numbers and let $Z_1,\dots,Z_n$ be sampled uniformly without replacement from $\{z_1,\dots,z_N\}$. Let
$$
\mu_N = \frac{1}{N}\sum_{i=1}^N z_i.
$$
Then for every $\varepsilon > 0$,
$$
\PP\!\left(
\left|
\frac{1}{n}\sum_{i=1}^n Z_i - \mu_N
\right| \ge \varepsilon
\right)
\le
2e^{-2n\varepsilon^2}.
$$
\end{theorem}

This inequality looks formally identical to the i.i.d.\ bounded case. The striking point is that sampling without replacement is \emph{not} independent, yet the same sub-Gaussian rate survives. The key structural reason is negative dependence: once a large value is sampled, it becomes unavailable later.

\subsubsection{McDiarmid's Inequality}

\begin{theorem}[McDiarmid inequality \cite{McDiarmid1989}]
\label{thm:mcdiarmid}
Let $X_1,\dots,X_n$ be independent random variables taking values in sets $\mathcal{X}_1,\dots,\mathcal{X}_n$. Let
$$
f:\mathcal{X}_1 \times \cdots \times \mathcal{X}_n \to \mathbb{R}
$$
be such that for every $i$ and every two inputs differing only in the $i$th coordinate,
$$
|f(x_1,\dots,x_n)-f(x_1,\dots,x_i',\dots,x_n)| \le c_i.
$$
Then for every $t>0$,
$$
\PP\!\left(
f(X_1,\dots,X_n)-\EE[f(X_1,\dots,X_n)] \ge t
\right)
\le
\exp\!\left(
-\frac{2t^2}{\sum_{i=1}^n c_i^2}
\right).
$$
The same holds for the lower tail, hence
$$
\PP\!\left(
\left|f(X_1,\dots,X_n)-\EE[f(X_1,\dots,X_n)]\right| \ge t
\right)
\le
2\exp\!\left(
-\frac{2t^2}{\sum_{i=1}^n c_i^2}
\right).
$$
\end{theorem}

This theorem controls \emph{stable functions of independent inputs}. It is especially useful when $f$ is not a linear statistic. In learning theory, it is the natural tool for bounding the fluctuation of a supremum such as a uniform deviation.

\subsubsection{Bernstein's Inequality}

\begin{theorem}[Bernstein inequality for bounded summands \cite{Bennett1962}]
\label{thm:bernstein}
Let $X_1,\dots,X_n$ be independent random variables with
$$
\EE[X_i]=0
\qquad \text{and} \qquad
|X_i| \le b
$$
almost surely for every $i$. Let
$$
V_n = \sum_{i=1}^n \EE[X_i^2].
$$
Then for every $t>0$,
$$
\PP\!\left(
\sum_{i=1}^n X_i \ge t
\right)
\le
\exp\!\left(
-\frac{t^2}{2(V_n + bt/3)}
\right).
$$
Hence
$$
\PP\!\left(
\left|
\sum_{i=1}^n X_i
\right| \ge t
\right)
\le
2\exp\!\left(
-\frac{t^2}{2(V_n + bt/3)}
\right).
$$
\end{theorem}

Bernstein improves on Hoeffding when the variance is much smaller than the worst-case range. For moderate deviations, the dominant scale is $V_n$ rather than $nb^2$.

\subsubsection{Comparison of the Four Inequalities}

\begin{itemize}
\item Theorem~\ref{thm:hoeffding} assumes independence and bounded support. It gives a universal sub-Gaussian tail.
\item Theorem~\ref{thm:hoeffding-without-replacement} removes independence but exploits the special finite-population structure of sampling without replacement.
\item Theorem~\ref{thm:mcdiarmid} no longer asks for a sum. Instead it asks for bounded coordinate influence.
\item Theorem~\ref{thm:bernstein} adds a variance term. When $V_n \ll nb^2$, it is much sharper than Hoeffding.
\end{itemize}

The main proof-level distinction is that Bernstein uses \emph{more distributional information} than Hoeffding. Hoeffding only remembers the range of each summand; Bernstein remembers both the range and the variance.

\subsection{Proofs}

\subsubsection{A Reusable MGF Lemma}

\begin{lemma}[Hoeffding lemma \cite{Hoeffding1963}]
\label{lem:hoeffding-lemma}
If $X$ is a random variable satisfying $a \le X \le b$ almost surely, then for every $\lambda \in \mathbb{R}$,
$$
\EE[e^{\lambda (X-\EE[X])}]
\le
\exp\!\left(\frac{\lambda^2 (b-a)^2}{8}\right).
$$
\end{lemma}

\begin{proof}
Let $m=\EE[X]$ and write
$$
p = \frac{m-a}{b-a} \in [0,1].
$$
By convexity of the exponential function, for every $x \in [a,b]$,
$$
e^{\lambda x}
\le
\frac{b-x}{b-a} e^{\lambda a}
+
\frac{x-a}{b-a} e^{\lambda b}.
$$
Taking expectations yields
$$
\EE[e^{\lambda X}]
\le
(1-p)e^{\lambda a} + pe^{\lambda b}.
$$
Therefore
$$
\EE[e^{\lambda (X-m)}]
\le
e^{-\lambda m}\bigl((1-p)e^{\lambda a} + pe^{\lambda b}\bigr).
$$
Set $h=\lambda(b-a)$. Since $m=a+p(b-a)$, the right-hand side becomes
$$
e^{-ph}(1-p+pe^h).
$$
Define
$$
g(h) = \log(1-p+pe^h) - ph.
$$
Then $g(0)=0$ and $g'(0)=0$. A direct computation gives
$$
g''(h)
=
\frac{pe^h(1-p)}{(1-p+pe^h)^2}
\le
\frac{1}{4}.
$$
Hence Taylor's theorem around $0$ implies
$$
g(h) \le \frac{h^2}{8}.
$$
Exponentiating,
$$
\EE[e^{\lambda (X-m)}]
\le
e^{h^2/8}
=
\exp\!\left(\frac{\lambda^2(b-a)^2}{8}\right).
$$
\end{proof}

\subsubsection{Proof of Hoeffding's Inequality}

\begin{proof}[Proof of Theorem~\ref{thm:hoeffding}]
Fix $\lambda >0$. By independence,
$$
\EE[e^{\lambda S_n}]
=
\prod_{i=1}^n \EE[e^{\lambda (X_i-\mu_i)}].
$$
Applying Lemma~\ref{lem:hoeffding-lemma} to each term,
$$
\EE[e^{\lambda S_n}]
\le
\exp\!\left(
\frac{\lambda^2}{8}
\sum_{i=1}^n (b_i-a_i)^2
\right).
$$
Now Lemma~\ref{lem:chernoff-template} gives
$$
\PP(S_n \ge t)
\le
\exp\!\left(
-\lambda t + \frac{\lambda^2}{8}\sum_{i=1}^n (b_i-a_i)^2
\right).
$$
The right-hand side is minimized at
$$
\lambda
=
\frac{4t}{\sum_{i=1}^n (b_i-a_i)^2}.
$$
Substituting this value gives
$$
\PP(S_n \ge t)
\le
\exp\!\left(
-\frac{2t^2}{\sum_{i=1}^n (b_i-a_i)^2}
\right).
$$
The lower-tail bound follows by applying the same argument to $-X_i$, and the two-sided bound follows by the union bound.
\end{proof}

\subsubsection{The Key Trick for Sampling Without Replacement}

\begin{theorem}[Hoeffding's convex comparison principle \cite{Hoeffding1963}]
\label{thm:hoeffding-comparison}
Let $z_1,\dots,z_N$ be fixed real numbers. Let $Z_1,\dots,Z_n$ be sampled without replacement from this population, and let $W_1,\dots,W_n$ be sampled with replacement from the same population. Then for every convex function $\varphi:\mathbb{R}\to\mathbb{R}$,
$$
\EE\!\left[\varphi\!\left(\sum_{i=1}^n Z_i\right)\right]
\le
\EE\!\left[\varphi\!\left(\sum_{i=1}^n W_i\right)\right].
$$
\end{theorem}

I am using Theorem~\ref{thm:hoeffding-comparison} as the imported comparison theorem from Hoeffding's paper. It is the key proof trick for Theorem~\ref{thm:hoeffding-without-replacement}: without-replacement sampling is more concentrated than with-replacement sampling for every convex test function.

\begin{proof}[Proof of Theorem~\ref{thm:hoeffding-without-replacement}]
Let
$$
\mu_N = \frac{1}{N}\sum_{i=1}^N z_i.
$$
Let $W_1,\dots,W_n$ be sampled with replacement from the same population. Then the $W_i$ are i.i.d., each lies in $[0,1]$, and each has mean $\mu_N$.

Apply Theorem~\ref{thm:hoeffding-comparison} with the convex function
$$
\varphi(x)=e^{\lambda x}.
$$
This gives, for every $\lambda>0$,
$$
\EE\!\left[
e^{\lambda \sum_{i=1}^n (Z_i-\mu_N)}
\right]
\le
\EE\!\left[
e^{\lambda \sum_{i=1}^n (W_i-\mu_N)}
\right].
$$
By Theorem~\ref{thm:hoeffding} applied to the i.i.d.\ variables $W_i \in [0,1]$,
$$
\PP\!\left(
\frac{1}{n}\sum_{i=1}^n W_i - \mu_N \ge \varepsilon
\right)
\le
e^{-2n\varepsilon^2}.
$$
Since the without-replacement sample has no larger exponential moments than the with-replacement sample, the same Chernoff optimization gives
$$
\PP\!\left(
\frac{1}{n}\sum_{i=1}^n Z_i - \mu_N \ge \varepsilon
\right)
\le
e^{-2n\varepsilon^2}.
$$
Apply the same argument to $1-Z_i$ to control the lower tail, then use the union bound to obtain
$$
\PP\!\left(
\left|
\frac{1}{n}\sum_{i=1}^n Z_i - \mu_N
\right|
\ge
\varepsilon
\right)
\le
2e^{-2n\varepsilon^2}.
$$
\end{proof}

\subsubsection{Proof of McDiarmid's Inequality}

\begin{proof}[Proof of Theorem~\ref{thm:mcdiarmid}]
Let
$$
M_i = \EE[f(X_1,\dots,X_n)\mid X_1,\dots,X_i],
\qquad i=0,\dots,n,
$$
with $M_0=\EE[f(X_1,\dots,X_n)]$ and $M_n=f(X_1,\dots,X_n)$. Then $(M_i)_{i=0}^n$ is the Doob martingale of $f$.

Define the martingale differences
$$
D_i = M_i - M_{i-1}.
$$
Then
$$
f(X_1,\dots,X_n)-\EE[f(X_1,\dots,X_n)]
=
\sum_{i=1}^n D_i.
$$

Fix $i$ and condition on $X_1,\dots,X_{i-1}$. If we vary only the $i$th coordinate while averaging over the future coordinates $X_{i+1},\dots,X_n$, the value of the conditional expectation can change by at most $c_i$. Therefore, conditional on $X_1,\dots,X_{i-1}$, the random variable $D_i$ has mean zero and lies in an interval of length at most $c_i$.

So the conditional version of Lemma~\ref{lem:hoeffding-lemma} gives
$$
\EE[e^{\lambda D_i}\mid X_1,\dots,X_{i-1}]
\le
\exp\!\left(\frac{\lambda^2 c_i^2}{8}\right).
$$
Iterating this bound,
$$
\EE\!\left[
e^{\lambda \sum_{i=1}^n D_i}
\right]
\le
\exp\!\left(
\frac{\lambda^2}{8}\sum_{i=1}^n c_i^2
\right).
$$
Applying Lemma~\ref{lem:chernoff-template},
$$
\PP\!\left(
\sum_{i=1}^n D_i \ge t
\right)
\le
\exp\!\left(
-\lambda t + \frac{\lambda^2}{8}\sum_{i=1}^n c_i^2
\right).
$$
Optimizing at
$$
\lambda
=
\frac{4t}{\sum_{i=1}^n c_i^2}
$$
gives
$$
\PP\!\left(
f(X_1,\dots,X_n)-\EE[f(X_1,\dots,X_n)] \ge t
\right)
\le
\exp\!\left(
-\frac{2t^2}{\sum_{i=1}^n c_i^2}
\right).
$$
The lower-tail and two-sided bounds follow as before.
\end{proof}

\subsubsection{A Scalar Inequality for Bernstein's Proof}

\begin{lemma}
\label{lem:bernstein-scalar}
For every real $u<3$,
$$
e^u \le 1 + u + \frac{u^2}{2(1-u/3)}.
$$
\end{lemma}

\begin{proof}
Consider
$$
q(u) = \log\!\left(1 + u + \frac{u^2}{2(1-u/3)}\right) - u.
$$
A direct calculation shows $q(0)=0$, $q'(0)=0$, and $q'(u)\ge 0$ for $u<3$. Hence $q(u)\ge 0$, which is equivalent to the claimed inequality.
\end{proof}

\subsubsection{Proof of Bernstein's Inequality}

\begin{proof}[Proof of Theorem~\ref{thm:bernstein}]
Fix $\lambda \in (0,3/b)$. Since $|X_i|\le b$, we have $\lambda X_i < 3$ almost surely. Applying Lemma~\ref{lem:bernstein-scalar} with $u=\lambda X_i$ and using $\EE[X_i]=0$,
$$
\EE[e^{\lambda X_i}]
\le
1 + \frac{\lambda^2 \EE[X_i^2]}{2(1-\lambda b/3)}.
$$
Using $1+x \le e^x$,
$$
\EE[e^{\lambda X_i}]
\le
\exp\!\left(
\frac{\lambda^2 \EE[X_i^2]}{2(1-\lambda b/3)}
\right).
$$

By independence,
$$
\EE\!\left[e^{\lambda \sum_{i=1}^n X_i}\right]
\le
\exp\!\left(
\frac{\lambda^2 V_n}{2(1-\lambda b/3)}
\right).
$$
Chernoff's bound therefore yields
$$
\PP\!\left(
\sum_{i=1}^n X_i \ge t
\right)
\le
\exp\!\left(
-\lambda t + \frac{\lambda^2 V_n}{2(1-\lambda b/3)}
\right).
$$
Now choose
$$
\lambda = \frac{t}{V_n + bt/3}.
$$
Then $\lambda < 3/b$, so this choice is admissible. A direct substitution gives
$$
\PP\!\left(
\sum_{i=1}^n X_i \ge t
\right)
\le
\exp\!\left(
-\frac{t^2}{2(V_n + bt/3)}
\right).
$$
Applying the same argument to $-X_i$ gives the lower tail and hence the two-sided version.
\end{proof}

\subsubsection{Proof Tricks Worth Reusing}

\begin{itemize}
\item \textbf{Hoeffding:} reduce everything to an exponential-moment bound for a single bounded centered variable.
\item \textbf{Sampling without replacement:} compare to a more convenient model. The whole theorem is driven by the convex comparison principle, Theorem~\ref{thm:hoeffding-comparison}.
\item \textbf{McDiarmid:} convert a complicated function into a sum of martingale increments whose ranges are controlled one coordinate at a time.
\item \textbf{Bernstein:} pay attention to the second moment. The variance term is the extra structure that buys a sharper tail.
\end{itemize}

\subsection{Connection to the Week 3 ERM Proof}

\subsubsection{The Learning-Theoretic Goal}

Let
$$
\widehat{h} \in \arg\min_{h\in\cH} L_n(h)
$$
be an empirical risk minimizer, where
$$
L_n(h) = \frac{1}{n}\sum_{i=1}^n \ind[h(x_i)\ne y_i]
$$
and
$$
L_{\cD}(h) = \PP_{(x,y)\sim \cD}(h(x)\ne y).
$$
The goal is to prove a high-probability bound of the form
$$
L_{\cD}(\widehat{h})
\le
\inf_{h\in\cH} L_{\cD}(h)
+ O\!\left(
\sqrt{\frac{\log \Gamma_{\cH}(2n) + \log(1/\delta)}{n}}
\right).
$$

\subsubsection{Step 1: A Uniform Deviation Quantity}

Define
$$
\Delta(S) = \sup_{h\in\cH} |L_{\cD}(h)-L_n(h)|.
$$
If we can show that $\Delta(S)$ is small with high probability, then the usual ERM three-line comparison finishes the argument.

\begin{lemma}
\label{lem:erm-three-line}
On the event $\{\Delta(S)\le \varepsilon\}$,
$$
L_{\cD}(\widehat{h})
\le
\inf_{h\in\cH} L_{\cD}(h) + 2\varepsilon.
$$
\end{lemma}

\begin{proof}
Let $h^* \in \arg\min_{h\in\cH} L_{\cD}(h)$. Then
$$
L_{\cD}(\widehat{h})
\le
L_n(\widehat{h}) + \varepsilon
\le
L_n(h^*) + \varepsilon
\le
L_{\cD}(h^*) + 2\varepsilon
=
\inf_{h\in\cH} L_{\cD}(h) + 2\varepsilon.
$$
\end{proof}

\subsubsection{Step 2: McDiarmid Controls the Fluctuation Around the Mean}

\begin{lemma}
\label{lem:mcdiarmid-for-delta}
For every $\delta \in (0,1)$, with probability at least $1-\delta$ over $S\sim \cD^n$,
$$
\Delta(S)
\le
\EE[\Delta(S)] + \sqrt{\frac{\log(1/\delta)}{2n}}.
$$
\end{lemma}

\begin{proof}
Changing one sample point can change each empirical risk $L_n(h)$ by at most $1/n$, while $L_{\cD}(h)$ stays fixed. Therefore changing one coordinate of $S$ changes $\Delta(S)$ by at most $1/n$. Applying Theorem~\ref{thm:mcdiarmid} with $c_i=1/n$ for all $i$ gives
$$
\PP\!\left(
\Delta(S)-\EE[\Delta(S)] \ge t
\right)
\le
e^{-2nt^2}.
$$
Setting
$$
t = \sqrt{\frac{\log(1/\delta)}{2n}}
$$
proves the claim.
\end{proof}

\subsubsection{Step 3: Symmetrization}

Let $S' \sim \cD^n$ be an independent ghost sample, and let
$$
L_n'(h) = \frac{1}{n}\sum_{i=1}^n \ind[h(x_i')\ne y_i'].
$$

\begin{lemma}[Symmetrization]
\label{lem:symmetrization}
$$
\EE[\Delta(S)]
\le
\EE_{S,S'}\!\left[
\sup_{h\in\cH} |L_n'(h)-L_n(h)|
\right].
$$
\end{lemma}

\begin{proof}
For each fixed $h$,
$$
L_{\cD}(h)-L_n(h) = \EE_{S'}[L_n'(h)-L_n(h)].
$$
Therefore
$$
\Delta(S)
=
\sup_{h\in\cH} \left| \EE_{S'}[L_n'(h)-L_n(h)] \right|.
$$
By Jensen's inequality for the convex function $x \mapsto |x|$,
$$
\left| \EE_{S'}[L_n'(h)-L_n(h)] \right|
\le
\EE_{S'}\!\left[ |L_n'(h)-L_n(h)| \right].
$$
Taking the supremum over $h$ and then using that $\sup \EE \le \EE \sup$,
$$
\Delta(S)
\le
\EE_{S'}\!\left[
\sup_{h\in\cH} |L_n'(h)-L_n(h)|
\right].
$$
Now take expectation over $S$.
\end{proof}

\subsubsection{Step 4: Reduce to the Transductive Concentration Problem}

Condition on the combined $2n$ sample points appearing in $S$ and $S'$. After conditioning, the pair $(S,S')$ is just a random split of a fixed pool of size $2n$ into two groups of size $n$ each. For a fixed hypothesis $h$, define the pool labels
$$
z_j(h) = \ind[h(\widetilde{x}_j)\ne \widetilde{y}_j] \in \{0,1\},
\qquad j=1,\dots,2n,
$$
where $(\widetilde{x}_j,\widetilde{y}_j)$ runs over the conditioned pool.

Then
$$
L_n(h)
$$
and
$$
L_n'(h)
$$
are the averages of the same $2n$ fixed numbers over a random split into two parts of size $n$.

\begin{lemma}
\label{lem:conditional-tail}
For every $t>0$ and every conditioned $2n$-point pool,
$$
\PP\!\left(
\sup_{h\in\cH} |L_n'(h)-L_n(h)| > t
\ \middle|\
S \cup S'
\right)
\le
2\Gamma_{\cH}(2n)e^{-nt^2/2}.
$$
\end{lemma}

\begin{proof}
Fix $h$. Let
$$
L_{\mathrm{pool}}(h) = \frac{1}{2n}\sum_{j=1}^{2n} z_j(h).
$$
Since the two halves have the same size,
$$
L_n'(h)-L_n(h)
=
2\bigl(L_{\mathrm{pool}}(h)-L_n(h)\bigr).
$$
By Theorem~\ref{thm:hoeffding-without-replacement},
$$
\PP\!\left(
|L_n(h)-L_{\mathrm{pool}}(h)| > \frac{t}{2}
\ \middle|\
S \cup S'
\right)
\le
2e^{-nt^2/2}.
$$
Equivalently,
$$
\PP\!\left(
|L_n'(h)-L_n(h)| > t
\ \middle|\
S \cup S'
\right)
\le
2e^{-nt^2/2}.
$$

Now, on the conditioned pool, two hypotheses that induce the same binary labeling on the $2n$ points have exactly the same value of $|L_n'(h)-L_n(h)|$. So there are at most
$$
\Gamma_{\cH}(2n)
$$
different quantities to consider. A union bound gives the claim.
\end{proof}

\subsubsection{Step 5: Tail Integration Bounds the Mean}

\begin{lemma}
\label{lem:expected-delta-bound}
$$
\EE[\Delta(S)]
\le
2\sqrt{
\frac{2(\log(2\Gamma_{\cH}(2n))+1)}{n}
}.
$$
\end{lemma}

\begin{proof}
By Lemma~\ref{lem:symmetrization},
$$
\EE[\Delta(S)]
\le
\EE\!\left[
\sup_{h\in\cH} |L_n'(h)-L_n(h)|
\right].
$$
For any nonnegative random variable $Y$,
$$
\EE[Y] = \int_0^\infty \PP(Y>t)\,dt.
$$
Applying this identity with
$$
Y = \sup_{h\in\cH} |L_n'(h)-L_n(h)|
$$
and using Lemma~\ref{lem:conditional-tail} then removing the conditioning,
$$
\EE[Y]
\le
\int_0^\infty \min\!\left\{1, 2\Gamma_{\cH}(2n)e^{-nt^2/2}\right\}\,dt.
$$

Set
$$
a = \sqrt{\frac{2(\log(2\Gamma_{\cH}(2n))+1)}{n}}.
$$
Then
$$
2\Gamma_{\cH}(2n)e^{-na^2/2} = e^{-1}.
$$
Hence
$$
\EE[Y]
\le
a + \int_a^\infty 2\Gamma_{\cH}(2n)e^{-nt^2/2}\,dt.
$$
Using $t^2-a^2 \ge 2a(t-a)$ for $t\ge a$,
$$
\int_a^\infty 2\Gamma_{\cH}(2n)e^{-nt^2/2}\,dt
\le
2\Gamma_{\cH}(2n)e^{-na^2/2}\int_a^\infty e^{-na(t-a)}\,dt
=
\frac{1}{ena}.
$$
Since $\log(2\Gamma_{\cH}(2n))+1 \ge 1$, we have
$$
\frac{1}{ena} \le a.
$$
Therefore
$$
\EE[Y] \le 2a
=
2\sqrt{
\frac{2(\log(2\Gamma_{\cH}(2n))+1)}{n}
}.
$$
\end{proof}

\subsubsection{Step 6: The Final ERM Guarantee}

\begin{theorem}
\label{thm:hw3-growth-erm}
Let
$$
\widehat{h} \in \arg\min_{h\in\cH} L_n(h).
$$
Then for every $\delta \in (0,1)$, with probability at least $1-\delta$ over $S\sim \cD^n$,
$$
L_{\cD}(\widehat{h})
\le
\inf_{h\in\cH} L_{\cD}(h)
+ 4\sqrt{
\frac{2(\log(2\Gamma_{\cH}(2n))+1)}{n}
}
+ 2\sqrt{\frac{\log(1/\delta)}{2n}}.
$$
In particular,
$$
L_{\cD}(\widehat{h})
\le
\inf_{h\in\cH} L_{\cD}(h)
+ O\!\left(
\sqrt{\frac{\log \Gamma_{\cH}(2n)+\log(1/\delta)}{n}}
\right).
$$
\end{theorem}

\begin{proof}
By Lemma~\ref{lem:mcdiarmid-for-delta} and Lemma~\ref{lem:expected-delta-bound}, with probability at least $1-\delta$,
$$
\Delta(S)
\le
2\sqrt{
\frac{2(\log(2\Gamma_{\cH}(2n))+1)}{n}
}
+ \sqrt{\frac{\log(1/\delta)}{2n}}.
$$
Apply Lemma~\ref{lem:erm-three-line} with this value of $\varepsilon$.
\end{proof}

\subsubsection{What Each Inequality Contributed}

\begin{itemize}
\item Theorem~\ref{thm:hoeffding} proves concentration for independent bounded averages.
\item Theorem~\ref{thm:hoeffding-without-replacement} is the transductive concentration step after conditioning on the $2n$ sampled points.
\item Theorem~\ref{thm:mcdiarmid} controls the fluctuation of the supremum $\Delta(S)$ around its mean.
\item Theorem~\ref{thm:bernstein} is not needed for the Week 3 ERM theorem as stated, but it is the next natural refinement when variance information is available and one wants sharper rates than Hoeffding can provide.
\end{itemize}

\subsection{References Used}

\begin{thebibliography}{9}

\bibitem{Hoeffding1963}
Wassily Hoeffding.
\newblock Probability Inequalities for Sums of Bounded Random Variables.
\newblock \emph{Journal of the American Statistical Association}, 58(301):13--30, 1963.
\newblock Available via Taylor \& Francis:
\url{https://www.tandfonline.com/doi/abs/10.1080/01621459.1963.10500830}.

\bibitem{Serfling1974}
Robert J. Serfling.
\newblock Probability Inequalities for the Sum in Sampling Without Replacement.
\newblock \emph{The Annals of Statistics}, 2(1):39--48, 1974.
\newblock Record:
\url{https://cir.nii.ac.jp/crid/1570854175574862336}.

\bibitem{McDiarmid1989}
Colin McDiarmid.
\newblock On the Method of Bounded Differences.
\newblock In \emph{Surveys in Combinatorics}, pages 148--188. Cambridge University Press, 1989.
\newblock Publisher page:
\url{https://www.cambridge.org/core/books/surveys-in-combinatorics-1989/on-the-method-of-bounded-differences/BD71F4D1C287C6E7450A7E12D4C3E082}.

\bibitem{Bennett1962}
George Bennett.
\newblock Probability Inequalities for the Sum of Independent Random Variables.
\newblock \emph{Journal of the American Statistical Association}, 57(297):33--45, 1962.
\newblock Available via Taylor \& Francis:
\url{https://www.tandfonline.com/doi/abs/10.1080/01621459.1962.10482149}.

\end{thebibliography}

\newpage
\section{Part B: The No-Free-Lunch Theorem and the Fundamental Theorem}

\subsection{Problem 1}

\subsubsection{Quantifiers and Adversarial Choices}

The theorem quantifies over:
\begin{itemize}
\item every learning algorithm $A$,
\item every sample size $n$ satisfying $n<|X|/2$.
\end{itemize}

After these are fixed, the adversary is allowed to choose:
\begin{itemize}
\item a set of $2n$ distinct domain points,
\item a bad target labeling on those points,
\item the realizable distribution supported on those labeled points.
\end{itemize}

The proof first averages over all labelings of a fixed $2n$-point subset, and then uses the probabilistic method to extract one specific bad labeling.

\subsubsection{The Averaging Construction}

Fix any learning algorithm $A$ and any $n$ with $n<|X|/2$. Since $|X|>2n$, we may choose distinct points
$$
x_1,\dots,x_{2n} \in X.
$$

For each labeling vector
$$
b=(b_1,\dots,b_{2n}) \in \{0,1\}^{2n},
$$
define a target function $f_b:X\to\{0,1\}$ by
$$
f_b(x_i)=b_i
\qquad \text{for } i=1,\dots,2n,
$$
and define it arbitrarily outside $\{x_1,\dots,x_{2n}\}$.

Now define a distribution $D_b$ on $X\times\{0,1\}$ by
$$
D_b\bigl((x_i,b_i)\bigr)=\frac{1}{2n},
\qquad i=1,\dots,2n.
$$
Thus $D_b$ is uniform on the $2n$ labeled points, and
$$
L_{D_b}(f_b)=0.
$$
So every $D_b$ is realizable.

\subsubsection{Average Population Loss Over Random Labelings}

Let $B=(B_1,\dots,B_{2n})$ be uniformly distributed on $\{0,1\}^{2n}$. An equivalent sampling procedure for $S\sim D_B^n$ is:
\begin{enumerate}
\item sample $B$ uniformly from $\{0,1\}^{2n}$,
\item sample indices $J_1,\dots,J_n$ i.i.d.\ uniformly from $[2n]$, independently of $B$,
\item output
$$
S=\bigl((x_{J_1},B_{J_1}),\dots,(x_{J_n},B_{J_n})\bigr).
$$
\end{enumerate}

Let
$$
V(S)=\{i\in[2n]: x_i \text{ appears at least once in } S\}
$$
be the set of observed points, and let
$$
U(S)=[2n]\setminus V(S)
$$
be the unseen points. Since the sample contains only $n$ examples,
$$
|V(S)| \le n,
$$
so
$$
|U(S)| \ge n.
$$

Now condition on the realized sample $S$. The hypothesis
$$
\widehat{h}=A(S)
$$
is then fixed. For any $i\in U(S)$, the bit $B_i$ was never revealed in the sample, so conditional on $S$ it is still an unbiased fair coin, independent of $\widehat{h}(x_i)$. Therefore
$$
\EE\!\left[\ind[\widehat{h}(x_i)\ne B_i]\mid S\right]=\frac{1}{2}
\qquad \text{for every } i\in U(S).
$$

Hence
$$
\EE\!\left[L_{D_B}(A(S))\mid S\right]
=
\frac{1}{2n}\sum_{i=1}^{2n}
\EE\!\left[\ind[\widehat{h}(x_i)\ne B_i]\mid S\right]
\ge
\frac{|U(S)|}{2n}\cdot \frac{1}{2}
\ge
\frac{1}{4}.
$$
Taking expectation over $S$ and $B$ gives
$$
\EE_{B}\EE_{S\sim D_B^n}[L_{D_B}(A(S))] \ge \frac{1}{4}.
$$

\subsubsection{Extracting a Single Bad Labeling}

Since the average over all $b\in\{0,1\}^{2n}$ is at least $1/4$, there exists at least one labeling $b^*$ such that
$$
\EE_{S\sim D_{b^*}^n}[L_{D_{b^*}}(A(S))] \ge \frac{1}{4}.
$$
Let
$$
f^*=f_{b^*},
\qquad
D=D_{b^*}.
$$
Then $L_D(f^*)=0$, so the distribution is realizable.

\subsubsection{From Expected Loss to Constant-Probability Failure}

Let
$$
Z = L_D(A(S)) \in [0,1].
$$
We already know
$$
\EE[Z] \ge \frac{1}{4}.
$$
Set
$$
p = \PP(Z \ge 1/8).
$$
Since $0\le Z\le 1$, we have
$$
\EE[Z]
\le
(1-p)\cdot \frac{1}{8} + p\cdot 1
=
\frac{1}{8} + \frac{7p}{8}.
$$
Combining this with $\EE[Z]\ge 1/4$ yields
$$
\frac{1}{4}
\le
\frac{1}{8}+\frac{7p}{8},
$$
so
$$
p \ge \frac{1}{7}.
$$

Therefore
$$
\PP_{S\sim D^n}\bigl(L_D(A(S))\ge 1/8\bigr)\ge \frac{1}{7}.
$$
This proves the No-Free-Lunch theorem.

\subsection{Problem 2}

\subsubsection{Corollary Closing the Lower-Bound Direction}

\begin{corollary}
\label{cor:vc-infinite-not-pac}
If $\operatorname{VCdim}(\cH)=\infty$, then $\cH$ is not realizably PAC learnable. Consequently, $\cH$ is not agnostically PAC learnable either.
\end{corollary}

\begin{proof}
Suppose $\operatorname{VCdim}(\cH)=\infty$. Fix any learning algorithm $A$ and any sample size $n$.

Because the VC dimension is infinite, there exists a set
$$
C=\{x_1,\dots,x_{2n}\}\subseteq X
$$
that is shattered by $\cH$. Apply the No-Free-Lunch theorem from Problem 1 to these $2n$ points. We obtain a distribution $D$ supported on $C$ and a labeling function
$$
f^*:C\to\{0,1\}
$$
such that
$$
\PP_{S\sim D^n}\bigl(L_D(A(S))\ge 1/8\bigr)\ge \frac{1}{7}.
$$

Since $C$ is shattered by $\cH$, there exists some $h^*\in\cH$ satisfying
$$
h^*(x)=f^*(x)
\qquad \text{for all } x\in C.
$$
Because $D$ is supported on $C$, this implies
$$
L_D(h^*)=0.
$$
So for every sample size $n$ there exists a realizable distribution $D$ on which algorithm $A$ still has failure probability at least $1/7$ at target excess level $1/8$.

Taking
$$
\varepsilon = \frac{1}{8},
\qquad
\delta = \frac{1}{7},
$$
we see that no sample size can satisfy the realizable PAC requirement uniformly over all distributions. Hence $\cH$ is not realizably PAC learnable.

Agnostic PAC learnability is stronger than realizable PAC learnability, so agnostic PAC learnability also fails.
\end{proof}

\subsubsection{Concrete Application to a Class of Infinite VC Dimension}

Consider the class
$$
\cH_{\mathrm{fin}}
=
\left\{
x \mapsto \ind[x\in F] : F\subseteq \mathbb{N} \text{ is finite}
\right\}.
$$

I claim that
$$
\operatorname{VCdim}(\cH_{\mathrm{fin}})=\infty.
$$

Indeed, let
$$
C=\{x_1,\dots,x_m\}\subseteq \mathbb{N}
$$
be any finite set. For any labeling
$$
y_1,\dots,y_m \in \{0,1\},
$$
define
$$
F=\{x_i : y_i=1\}.
$$
Then $F$ is finite, so the hypothesis
$$
h_F(x)=\ind[x\in F]
$$
belongs to $\cH_{\mathrm{fin}}$, and it satisfies
$$
h_F(x_i)=y_i
\qquad \text{for all } i=1,\dots,m.
$$
Thus every finite set is shattered, so the VC dimension is infinite.

Applying Corollary~\ref{cor:vc-infinite-not-pac}, $\cH_{\mathrm{fin}}$ is not PAC learnable. More explicitly: for every $n$ and every learning algorithm $A$, there exists a realizable distribution $D_n$ supported on $\{1,\dots,2n\}$ and a target hypothesis $h_n^*\in \cH_{\mathrm{fin}}$ such that
$$
\PP_{S\sim D_n^n}\bigl(L_{D_n}(A(S))\ge 1/8\bigr)\ge \frac{1}{7}.
$$
So this class gives a concrete example of the lower-bound direction of the Fundamental Theorem.

\subsection{Problem 3}

\subsubsection{A Concrete Learning Rule}

Fix a sample size $n$ and domain
$$
X=\{1,2,\dots,2n\}.
$$
Define the learning rule $A_{\mathrm{mem}}$ as follows: given a training sample
$$
S=((x_1,y_1),\dots,(x_n,y_n)),
$$
output the classifier
$$
A_{\mathrm{mem}}(S)(x)
=
\ind[\text{$x$ appears in $S$ with label $1$}].
$$
So the learner predicts $1$ exactly on points it has explicitly seen labeled $1$, and predicts $0$ everywhere else.

\subsubsection{Bad Distribution and Target}

Let the target function be
$$
f^*(x)=1
\qquad \text{for all } x\in X.
$$
Let $D$ be the uniform distribution on
$$
\{(1,1),(2,1),\dots,(2n,1)\}.
$$
Then
$$
L_D(f^*)=0.
$$
Moreover, $f^*$ belongs to $\cH_{\mathrm{fin}}$, since it is the indicator of the finite set $\{1,\dots,2n\}$.

\subsubsection{Failure Calculation}

Given any sample $S\sim D^n$, let
$$
V(S)=\{x\in X : x \text{ appears at least once in } S\}.
$$
Since the sample has only $n$ examples,
$$
|V(S)|\le n.
$$
Therefore the unseen set
$$
U(S)=X\setminus V(S)
$$
satisfies
$$
|U(S)|\ge n.
$$

For every $x\in U(S)$, the rule $A_{\mathrm{mem}}(S)$ predicts $0$, while the target label is $f^*(x)=1$. Hence
$$
L_D(A_{\mathrm{mem}}(S))
=
\frac{|U(S)|}{2n}
\ge
\frac{n}{2n}
=
\frac{1}{2}.
$$

So the learner fails not merely with constant probability, but deterministically:
$$
\PP_{S\sim D^n}\bigl(L_D(A_{\mathrm{mem}}(S))\ge 1/2\bigr)=1.
$$
In particular,
$$
\PP_{S\sim D^n}\bigl(L_D(A_{\mathrm{mem}}(S))\ge 1/8\bigr)=1.
$$

This is an explicit witness to the kind of failure that the abstract NFL theorem guarantees existentially.

\subsection{Problem 4}

\subsubsection{Learning Setting}

The Week 1 theorem lives in the \emph{online transductive} mistake-bound setting: the learner receives one point at a time and is judged by the total number of mistakes on a finite domain.

The Week 3 theorem lives in the \emph{batch i.i.d.\ PAC} setting: the learner receives an i.i.d.\ training sample and is judged by the population loss under an unknown distribution.

\subsubsection{Adversary Type}

In Week 1, the proof is driven by an \emph{adaptive adversary}. At each round, the adversary can look at the learner's current prediction on a fresh point and then choose the opposite label. After the interaction ends, this adaptive construction can be repackaged as a single bad function $f$ and a fixed enumeration of the domain.

In Week 3, the adversary is \emph{oblivious}. It chooses a bad distribution $D$ and a realizable target function $f^*$ before the sample is drawn. The learner is then evaluated on the randomness of the i.i.d.\ sample, not on an adaptively chosen online sequence.

\subsubsection{Form of the Conclusion}

The Week 1 conclusion is deterministic and combinatorial:
$$
\text{there exists a sequence on which the learner makes } |X| \text{ mistakes.}
$$
It is a worst-case online impossibility result.

The Week 3 conclusion is probabilistic and statistical:
$$
\PP_{S\sim D^n}\bigl(L_D(A(S))\ge 1/8\bigr)\ge \frac{1}{7}.
$$
It is a lower bound on generalization performance under random sampling.

\subsubsection{Why Both Are Called No-Free-Lunch}

Both theorems say the same qualitative thing: without structural restrictions on the target class, no learner can succeed uniformly over all problems.

Week 1 says there is no universally low-mistake online learner on an unrestricted finite domain.

Week 3 says there is no universally PAC learner on an unrestricted class of labelings, because for every sample size there is a realizable distribution on which the learner still has constant population error with constant probability.

So the two theorems live in different models and produce quantitatively different lower bounds, but they both express the same underlying message: without a complexity restriction, learning has no free lunch.

\newpage
\section{Part C: AI Usage Report}

\subsection*{1. Parts of the Assignment Where I Used AI}

This time it keeps the same, but here I added some new skills and help unify the notations. See details as next sections.

\subsection*{2. Concrete AI Suggestions I Accepted}

One suggestion I accepted was to separate each concentration inequality into the same template: theorem statement, assumptions, proof sketch, full proof, and proof trick.

Another suggestion I accepted was to present the i.i.d.\ ERM argument through the quantity
$$
\Delta(S)=\sup_{h\in\cH}|L_{\cD}(h)-L_n(h)|
$$
and then split the proof into McDiarmid for fluctuation, symmetrization for the mean, and transductive concentration after conditioning on $S\cup S'$.

\subsection*{3. Concrete AI Suggestions I Rejected or Substantially Modified}

I rejected an early proof for the without-replacement Hoeffding inequality that tried to treat the sample as if it were independent.

\subsection*{4. How I Verified Correctness}

For theorem statements in Part A, I checked the formulas and assumptions against the cited references and against the Week 3 slides. Also this part is covered in the previous course, so it's easy to check.

For the ERM proof, I checked each reduction separately: McDiarmid for the fluctuation term, symmetrization for replacing population risk by ghost-sample empirical risk, and the conditional transductive argument for the growth-function term $\Gamma_{\cH}(2n)$.

\subsection*{5. The 5 Most Recent AI Workflow Changes and Why}

\textbf{(i) Added the `theorm-ref` skill.} This came directly from this assignment. I needed a workflow for theorem-heavy writing that forces theorem statements, labels, source citations, unified notation, proof sketches, full proofs, and proof-trick summaries.

\textbf{(ii) Added a citation-first rule for theorem sections.} For concentration inequalities, I now want the source attached at the theorem statement stage instead of added at the end. This reduces the chance of drifting into uncited folklore formulations.

\textbf{(iii) Added a unified-notation rule for theorem collections.} In earlier drafts it was too easy for different theorems to use unrelated symbols. For this assignment I standardized notation for means, centered sums, empirical averages, variance terms, and bounded-difference constants.

\textbf{(iv) Added PDF-to-statement extraction with `pypdf`.} Instead of repeatedly rereading the raw PDF, I now extract the assignment into a reusable companion `.tex` file first. That made it easier to preserve the exact prompt and reuse it across turns.

\textbf{(v) Added a stronger compile-and-check loop for theory writeups.}

\end{document}
